% Figure and table captions for the ISC paper. Generated by analysis/paper_assets_isc.py; do not edit.
% Each macro holds the caption text after "Fig. N." (LaTeX adds the label).
\newcommand{\capFigForest}{Top-1 flips (a) and success flips (b) between pairs of platforms, with 95\% bootstrap intervals over targets. Bands give recompilation alone on the platform with the larger rate. FU is the follow-up. In FU1, success is a DockQ of at least 0.23.}
\newcommand{\capFigEcdf}{Cumulative distribution of the ligand RMSD between paired samples, for two pairs of platforms, recompilation on the L4, and a change of seed on the v6e. Recompilation on the v5e and on the v6e is bit-identical and not drawn. Values of 0 are drawn at $10^{-4}$~\AA{}. The curves use 1425 sample pairs on 57 targets, and 2850 for the change of seed. The three glycan targets have no ligand RMSD.}
\newcommand{\capFigMargin}{Exploratory, specified after the main results were read. Per-seed top-1 flip rate by quintile of the first condition's ranking margin, with 95\% bootstrap intervals over targets. Each quintile holds 60 pairs of predictions.}
\newcommand{\capFigPerturbation}{Exploratory. Top-1 flip rate of nine contrasts against the median shift in ranking score, divided by the median top-two margin. The dashed line is a two-sample model.}
\newcommand{\capFigPerturbationApp}{Exploratory. Top-1 flip rate of nine contrasts against the median shift in ranking score, divided by the median top-two margin. The dashed line is a two-sample model.}
\newcommand{\capFigSeeds}{Follow-up 8, outside the pre-registered study. Difference in the success rate of the overall top-ranked sample over the first n seeds, v5e minus L4, in points, with 95\% bootstrap intervals over 60 targets. The band is the $\pm$2 point criterion of Q3.}
\newcommand{\capFigCost}{Per structure, one target, one seed, five samples, fresh cache. At bucket 1024 the cost was 0.113 USD on demand, 0.0409 Spot and 0.0436 Flex-start on the v5e; 0.170, 0.102 and 0.0771 on the v6e; 0.173 and 0.104 on the L4; 0.0932 and 0.0530 on the L4 Triton arm. All runs were on demand: Spot and Flex-start figures are list prices read on 8 and 9 October 2026 applied to these times, excluding VM setup, idle time and boot disk. Google publishes no Flex-start price for g2-standard-8.}
\newcommand{\capFigSize}{Follow-up 10, outside the pre-registered study: for both primary pairs, top-1 flip minus recompilation showed no size difference detected at this sample size: large minus control, $-$1.8 [$-$10.6, 6.5] points for the v5e-v6e difference beyond recompilation on the v5e and +2.9 [$-$7.1, 12.9] points for the v6e-L4 difference beyond recompilation on the v6e (interval widths 17.1 and 20.0).}
\newcommand{\capTabRuns}{Runs per study}
\newcommand{\capTabSPertFive}{Exploratory (c) at five seeds}
\newcommand{\capTabArms}{Arms of the Main Study and of Follow-Ups 3 and 9}
\newcommand{\capTabHeadline}{Flips and Structure Differences in the Main Study}
\newcommand{\capTabQtwo}{Time, Cost, and Memory per Structure (Q2)}
\newcommand{\capTabSFuFifteenRC}{Follow-up 15: recompilation of the float32 trunk and of the L4 Triton}
\newcommand{\capTabSFuFifteenDF}{Follow-up 15: $D_F$ of the H100 and the A100 with recompilation measured}
\newcommand{\capTabSFuSixteen}{Follow-up 16: the compiled programs of the v5e and the v6e before the Pairformer, for inputs padded to 256 tokens}
\newcommand{\capTabSFuSeventeenCells}{Follow-up 17: the data type of the trunk and the matrix-product precision, v6e~vs~L4}
\newcommand{\capTabSFuSeventeenPairs}{Follow-up 17: the four accelerators with the float32 trunk at the default precision}
\newcommand{\capTabSFuEighteen}{Follow-up 18: the H100 and the A100 with Triton attention}
\newcommand{\capTabSFuNineteenRC}{Follow-up 19: float32 recompilation on 60 targets}
\newcommand{\capTabSFuNineteenDF}{Follow-up 19: $D_F$ of the H100 and the A100 with recompilation measured on 60 targets}
\newcommand{\capTabSFuTwenty}{Follow-up 20: the L4 without autotuning}
\newcommand{\capTabSFuTwentyOne}{Follow-up 21: recompilation on a new VM of the same GPU type}
\newcommand{\capTabFollowups}{Follow-ups 1 to 21}
\newcommand{\capTabSFuThreeQThree}{Follow-up 3, Q3: success rate over the first n seeds}
\newcommand{\capTabSFuThreeQThreeDiff}{Follow-up 3, Q3: differences in success rate}
\newcommand{\capTabSFuThreeMemory}{Follow-up 3, M7: device memory of the A100 arms}
\newcommand{\capTabRules}{Reporting rules}
\newcommand{\capTabSRulesSeeds}{Ranking rules: seed consensus}
\newcommand{\capTabSRulesTies}{Ranking rules: tie tolerance}
\newcommand{\capTabSRulesTieShare}{Ranking rules: share of ties and conflicts}
\newcommand{\capTabSRulesRounded}{Ranking rules: full-precision against rounded score}
\newcommand{\capTabSRulesModel}{Exploratory: observed flip rate against two models}
\newcommand{\capTabSCampChanged}{Campaign simulation: reported samples that change}
\newcommand{\capTabSCampPoints}{Campaign simulation: success-rate difference}
\newcommand{\capFigFuThirteen}{(a) Top-1 flip rate on the v6e against the size of a random perturbation of the trunk output, in units of $u = 2^{-8}$ (seeds 1 and 2). (b) Top-1 flip rates with the bfloat16 and the float32 trunk (seeds 1 to 5), on a square-root scale. All rates use 60 targets and 95\% intervals over targets.}
\newcommand{\capFigFuThirteenL}{Follow-up 13, layer-capture arm: the relative difference between two platforms after each Pairformer block, in units of $u = 2^{-8}$, on 10 pilot targets with seed 1. Depth 0 is the input of the Pairformer, after the input embedding, the recycling embedding, and the MSA module. Lines give the median and bands the range over targets. The pair representation was captured at even depths on a fixed token subset. Recompilation on the L4 has a median difference of 0 at every depth and is not drawn on the log scale.}
\newcommand{\capTabSFuThirteenChecks}{Follow-up 13: checks before any result}
\newcommand{\capTabSFuThirteenP}{Follow-up 13, perturbation arm: flip rates against size}
\newcommand{\capTabSFuThirteenEps}{Follow-up 13, perturbation arm: the size that reaches each platform rate}
\newcommand{\capTabSFuThirteenL}{Follow-up 13, layer capture: difference at the trunk output after each recycle}
\newcommand{\capTabSFuThirteenF}{Follow-up 13, float32 trunk arm: flip rates per contrast}
\newcommand{\capTabSFuThirteenDF}{Follow-up 13, float32 trunk arm: reading and secondary comparisons}
\newcommand{\capTabSFuThirteenCost}{Follow-up 13, float32 trunk arm: time and memory against the bfloat16 trunk}
\newcommand{\capTabSFuFourteenChecks}{Follow-up 14: checks before any result}
\newcommand{\capTabSFuFourteenAgree}{Follow-up 14: agreement of the float32 trunk across platforms}
\newcommand{\capTabSFuFourteenDF}{Follow-up 14: the float32 trunk against the bfloat16 trunk}
\newcommand{\capTabSFuFourteenSec}{Follow-up 14: float32 against bfloat16 trunk, every pair}
\newcommand{\capTabSFuFourteenF}{Follow-up 14: float32 trunk, recompilation, and the deterministic flag}
\newcommand{\capTabSFuFourteenSuccess}{Follow-up 14: success rate per arm}
\newcommand{\capTabSFuFourteenCost}{Follow-up 14: time and compiled memory against the reference arm}
\newcommand{\capTabSFuFourteenMemory}{Follow-up 14: compiled memory per bucket}
\newcommand{\capTabBeyond}{The Difference Beyond the Main Study}
\newcommand{\capTabRepro}{The Cost of Reproducibility per Platform}
\newcommand{\capTabSESummary}{Exploratory (e): flips that matter, by the criterion that separates the two top-ranked structures}
\newcommand{\capTabSReturned}{Ligand RMSD of the returned structure across platforms (post hoc)}
\newcommand{\capTabSRare}{Rare-event readings under BCa intervals (post hoc)}
\newcommand{\capTabSWinners}{The two returned poses compared directly (post hoc)}
\newcommand{\capTabSWinnersExt}{The two returned poses compared directly, further contrasts (post hoc)}
\newcommand{\capTabSPbFlips}{Which success criterion changes in a success flip (post hoc)}
\newcommand{\capTabSBestTwentyFive}{Success of the best of all samples, per platform (post hoc)}
\newcommand{\capTabSGtat}{The L4 and the A100, both with Triton attention (post hoc)}
\newcommand{\capTabCampaign}{The campaign by label and by device}
\newcommand{\capTabProvisioning}{Provisioning by arm and follow-up}
